% OWNER: sections/s4_3_guidelines.tex. Edit only this file.
\seclabel
\claim{\secnum{4.3}{sec:guidelines}\quad \code{plan()} rejects sector breaches; \code{guidelines()} reports the other limits}

\lead{The contest sets 4 investment guidelines, and the server enforces only the stock limit when an order is placed (\figref{guiderules}).} The server blocks an order that would take a stock past its stock limit (\code{weiLimit}): 40\% for Samsung Electronics and 15\% for every other stock. Sector and small-cap breaches are only counted, as violations in the \code{Portfolio/Violations} reply, and the feed's compliance record (\code{cmpl}) carries this week's turnover. Turnover has to reach 5\% every week; up to 3 weekly violations are allowed, and 4 or more exclude the portfolio from the ranking and the prizes. The limits come from the guideline page of the contest site.

\lead{The values behind the guidelines come from 5 smaller feed fields beside the 4 main ones and from the server's \code{Violations} reply (\figref{guiderules}).} Besides \code{bs}, \code{pos}, \code{ords} and \code{msecs}, a snapshot carries \code{bm}, \code{bs_updated}, \code{cmpl}, \code{noPosPrcs} and \code{tags}. The \code{state} module keeps them as they arrive and lets a delta overwrite any of them that is not null; \code{bs_updated}, present in every delta, is merged into \code{bs}, so the net assets (\code{na}) stay current (Section~\ref{sec:feed}). \code{tags} labels stocks by ticker, and \code{msecs} holds one record per GICS sector with a record \code{Id}, the sector code \code{sec} and the market weight \code{wei}. The \code{Violations} reply gives the server's own counts in 3 parts (\code{Current}, \code{Cumulative} and \code{Details}); it has no turnover field, so turnover comes only from \code{cmpl}.

% OWNER: figs/guiderules.tex. Edit only this file.
\fig[wide]{guiderules}{The server blocks only stock-limit breaches and counts the other 3 guidelines}[Contest guidelines, how each is enforced, the fields behind it and its value at 09:52 KST on 24 Nov 2026]{%
{\tabfont\begin{tabularx}{\linewidth}{>{\RaggedRight\arraybackslash}p{19mm} >{\RaggedRight\arraybackslash}p{33mm} >{\RaggedRight\arraybackslash}p{43mm} >{\RaggedRight\arraybackslash}X >{\RaggedRight\arraybackslash}p{40mm}}
\thd{Rule} & \thd{Limit} & \thd{How enforced} & \thd{Value at 09:52} & \thd{Source field}\\
Stock weight & Up to the stock's \code{weiLimit}: Samsung Electronics 40\%, others 15\% & Server blocks the order. Wrapper: \code{plan()} caps each target at \code{weiLimit}, or at 15\% for a stock not held (Samsung Electronics 40\%), by water-filling; a buy past the limit is refused before sending & \code{StkViolCnt} 0 & \code{pos}: \code{wei}, \code{weiLimit}; \code{Violations}: \code{StkViolCnt}, \code{StockViols}\\\rowsep
Sector weight & Up to max(2 $\times$ market sector weight, 10\%) & Server counts violations. Wrapper: \code{plan()} rejects the target set with \code{ValueError} & \code{SectViolCnt} 0; least headroom 2.01\,pp (\figref{sectors}) & \code{pos}: \code{sec}, \code{wei}; \code{msecs}: \code{sec}, \code{wei} (market weight, \%); \code{Violations}: \code{SectViolCnt}, \code{SectorViols}\\\rowsep
Small caps & Stocks under KRW 1 trillion market cap, up to 30\% in total & Server counts violations. Wrapper reports the total through \code{guidelines()} & 3.35\% in 1 stock (\code{A090460}); \code{Excess} 0 & \code{tags} (label \code{SmallCap} by ticker); \code{Violations}: \code{SmallCapViol} (\code{TotalWei}, \code{Limit}, \code{Excess})\\\rowsep
Weekly turnover & At least 5\% every week; 4 or more weekly violations exclude the portfolio & Server counts weekly violations in the feed. Wrapper reports them through \code{guidelines()} & \code{currWto} 4.8\%, \code{avgWto} 28.0\%, \code{wtoViols} 1 & \code{cmpl}: \code{currWto}, \code{avgWto}, \code{wtoViols}, \code{expWtoViols}\\
\end{tabularx}}}{Limits from the guideline page of the contest site; values from the live feed (\code{pos}, \code{msecs}, \code{tags}, \code{cmpl}) and the \code{Portfolio/Violations} reply. \code{Violations} returns \code{Current} and \code{Cumulative} counts (all 0 at 09:52; \code{Cumulative} adds \code{ViolDays}) and \code{Details}. Share weights use net assets \code{bs.na} (KRW 837,002,427 at 09:52), which \code{bs_updated} keeps current. \code{bm}, \code{noPosPrcs} and \code{cmpl.asr} are kept but not interpreted; the meaning of \code{expWtoViols} (1 at 09:52) has not been confirmed. pp = \mbox{percentage} points.}{Live feed and \code{Portfolio/Violations}, DART Quant Team}


\Needspace{4\baselineskip}\lead{\code{guidelines()} combines both sources into one reply, and its \code{ok} flag needs 0 current violations and 5\% turnover this week.} The reply lists the 5 stocks with the least headroom under \code{weiLimit}, the weight, limit and headroom of every sector held (\figref{sectors}), the small-cap total, the turnover values of \code{cmpl} and the server's \code{Violations} reply unchanged. \code{ok} is True only if the server counts 0 current violations and this week's turnover is at least 5\% (D25).

% OWNER: figs/sectors.tex. Edit only this file.
\fig{sectors}{Consumer Discretionary was nearest its sector limit, 2.01\,pp below}[Sector weights (\%) against their limits, headroom in pp, largest holding first; 09:52 KST on 24 Nov 2026]{%
{\tabfont\begin{tabularx}{\linewidth}{>{\RaggedRight\arraybackslash}X R{19mm} R{15mm} R{14mm} R{16mm}}
\thd{GICS sector} & \thd{Market weight} & \thd{Portfolio} & \thd{Limit} & \thd{Headroom}\\
45 Information Technology & 52.90 & 45.74 & 105.80 & 60.06\\\rowsep
20 Industrials & 17.70 & 30.17 & 35.40 & 5.23\\\rowsep
\hl 25 Consumer Discretionary & 5.75 & 9.49 & 11.50 & 2.01\\\rowsep
50 Communication Services & 5.00 or less & 6.79 & 10.00 & 3.21\\\rowsep
35 Health Care & 5.69 & 6.54 & 11.38 & 4.84\\
\end{tabularx}}}{Shaded row: the sector with the least headroom. Limit = max(2 $\times$ market weight, 10\%); headroom = limit $-$ portfolio weight. IT, the largest sector, had 60.06\,pp to spare. Market weights are derived from the limits; for Communication Services the 10\% floor applies, so its market weight is only bounded. The 5 sectors hold 98.73\% of the portfolio.}{\code{guidelines()} output, DART Quant Team}


\lead{This week's turnover was 4.8\% at 09:52 on Tuesday, below the 5\% minimum, so \code{ok} was False while the server counted 0 violations.} \code{cmpl.currWto} read 4.8\% and \code{cmpl.wtoViols} read 1, against the 3 weekly violations the rules allow. The wrapper reports these numbers; what to trade and when is left to the strategy (N2).

\lead{Keying the market sector weights by the \code{msecs} record \code{Id} put IT 35.74~pp over its limit, while the server's \code{SectViolCnt} read 0.} The first version of \code{guidelines()} used the record \code{Id} as the sector. Record \code{Id} 1 holds sector 10, so a lookup for sector 45 found no market weight, the limit fell to the 10\% floor and IT at 45.74\% appeared over it. Keyed by the sector code \code{sec}, the IT limit is 105.80\%. The server's count of 0 exposed the defect.

\lead{\code{plan()} checks every target set against the contest rules before it creates any order.} In order, it raises \code{ValueError} when the targets sum above \code{max_total} (100\%); caps each target stock at its \code{weiLimit}, moving the excess to the other stocks by water-filling; raises when a target names a stock outside the universe; and raises when a sector's total exceeds max(2 $\times$ market weight, 10\%), the check of \code{sector_violations()} in \code{orders} (D23). It then subtracts the remaining size of working orders and drops orders below one share's weight (Section~\ref{sec:oneshare}). A stock the portfolio does not hold has no \code{pos} record and therefore no \code{weiLimit} in the feed, so \code{plan()} caps it at the guideline value, 15\% (Samsung Electronics 40\%); the check before each order is sent applies the same caps and refuses a buy that would pass them. The sector of each stock comes from \code{Universe}, which the client loads once at start (2,648 stocks); the market weights come from \code{msecs} at each call.

\lead{The wrapper checks sectors itself because the server reports a sector breach only as a count after the weights have moved.} \code{plan()} rejects the target set and leaves the excess where it is. Spreading a sector's excess, as water-filling does for a stock, would lower every stock in that sector and raise all the others, and each changed stock costs one of the 80 daily orders. The error names the sector, its total and its limit, so the caller can adjust the targets (\figref{sectortest}).

% OWNER: figs/sectortest.tex. Edit only this file.
\fig{sectortest}{\code{plan()} rejected Industrials at +6\,pp and accepted +3\,pp}[Industrials weight in \% (GICS 20, limit 35.36\%) moved by shifting weight from Samsung Electronics to HD Hyundai Heavy Industries; 09:55 KST on 24 Nov 2026]{%
{\tabfont\begin{tabularx}{\linewidth}{>{\RaggedRight\arraybackslash}p{58mm} R{15mm} >{\RaggedRight\arraybackslash\leftskip=8pt}X}
\thd{Target change} & \thd{Industrials} & \thd{\code{plan()} result}\\
None (current weights) & 30.14 & 0 orders\\\rowsep
\hl Industrials +6\,pp, Samsung Electronics $-$6\,pp & 36.14 & \code{ValueError}: \texttt{20 36.14\% > 35.36\%}; no orders\\\rowsep
Industrials +3\,pp, Samsung Electronics $-$3\,pp & 33.14 & 2 orders: sell Samsung Electronics 3\%, then buy HD Hyundai Heavy Industries 3\%\\
\end{tabularx}}}{Shaded row: the rejected target set. Samsung Electronics is in sector 45, whose limit is 105.80\%, so only Industrials could break a limit. Nothing was sent.}{DART Quant Team}


\lead{The Industrials limit read 35.40\% at 09:52 and 35.36\% at 09:55, so \code{plan()} reads every sector limit from the feed at each call.} The market weights in \code{msecs} change about every 12\,s (Section~\ref{sec:feed}), and a limit stored once would drift from the server's. In the accepted set of \figref{sectortest}, the orders came out sells first, so the sale provides the cash for the purchase.

\lead{The small-cap rule is reported only, because the wrapper has no market capitalisation for a stock it does not hold.} \code{tags} labelled 1 stock, the small cap the portfolio holds, so a target naming a new small cap cannot be recognised from the feed, and market data stays outside the wrapper (N1). At 3.35\% against 30\%, the margin was 26.65~pp.
\verdict{\code{plan()} rejects targets above 100\%, outside the universe or over a sector limit read from the feed at each call, before any order exists (D19, D23). \code{guidelines()} reports the other limits with the server's \code{Violations} reply; its \code{ok} is True only with 0 current violations and at least 5\% turnover this week (D25).}
